\documentclass[preprint,12pt]{elsarticle}
\usepackage{lmodern}
\usepackage{amsmath,amssymb}
\usepackage{graphicx}
\usepackage{booktabs}
\usepackage{multirow}
\usepackage{array}
\usepackage{xcolor}
\usepackage[section]{placeins}
\usepackage{float}
\usepackage{xr-hyper}
\usepackage{hyperref}
\externaldocument[S-]{supplementary}
\hypersetup{hidelinks,
  pdftitle={Calibration-Free Post-Estimation Fusion of Dual-View 3D Human Pose Streams: Robustness, Front-End Generalization, and Multi-View Extension},
  pdfauthor={Kaixu Chen, Qingyu Wang, Nan Luo, Chun Xie, Eiichi Naito, Toshitaka Kimura, Itaru Kitahara}}
\usepackage{lineno}
\setlength{\emergencystretch}{2em}
\journal{Image and Vision Computing}

% Upright prose, including class-generated headings and bibliography emphasis.
\AddToHook{begindocument/end}{%
  \renewcommand{\itshape}{\fontshape{n}\selectfont}%
  \renewcommand{\slshape}{\fontshape{n}\selectfont}%
}
\renewcommand{\textit}[1]{\textup{#1}}
\renewcommand{\emph}[1]{\textup{#1}}

\begin{document}

\begin{frontmatter}

\title{Calibration-Free Post-Estimation Fusion of Dual-View 3D Human Pose Streams: Robustness, Front-End Generalization, and Multi-View Extension\tnoteref{tnote:conference}}
\tnotetext[tnote:conference]{A preliminary version of this work has been accepted for publication at the 9th International Workshop on Multimedia Content Analysis in Sports (MMSports 2026). This article substantially extends that work.}

\author[ccs]{Kaixu Chen\corref{cor1}}
\ead{chenkaixusan@gmail.com}
\author[imiss]{Qingyu Wang}
\ead{wang.qingyu@image.iit.tsukuba.ac.jp}
\author[imiss]{Nan Luo}
\ead{nan.luo@image.iit.tsukuba.ac.jp}
\author[ccs]{Chun Xie}
\ead{xie.chun@image.iit.tsukuba.ac.jp}
\author[nict]{Eiichi Naito}
\ead{eiichi.naito@nict.go.jp}
\author[ntt]{Toshitaka Kimura}
\ead{toshitaka.kimura@ntt.com}
\author[ccs]{Itaru Kitahara}
\ead{kitahara.itaru@image.iit.tsukuba.ac.jp}

\cortext[cor1]{Corresponding author.}
\address[ccs]{Center for Computational Sciences, University of Tsukuba, Tsukuba, Ibaraki, Japan}
\address[imiss]{Intelligent and Mechanical Interaction Systems, University of Tsukuba, Tsukuba, Ibaraki, Japan}
\address[nict]{Advanced ICT Research Institute, National Institute of Information and Communications Technology, Osaka, Japan}
\address[ntt]{NTT Communication Science Laboratories, Tokyo, Japan}

\begin{abstract}
Monocular three-dimensional (3D) human pose estimates remain vulnerable to depth ambiguity, self-occlusion, and temporal jitter. A second camera can provide complementary evidence, but exploiting it through calibrated reconstruction is difficult when cameras follow a distant, fast-moving athlete or are installed temporarily. We present CanonFuse3D, a calibration-free post-estimation framework that fuses two independently estimated, approximately synchronized monocular 3D pose streams without requiring images, 2D observations, or camera parameters at the fusion stage. Human-centric canonicalization and sequence-level alignment expose cross-view agreement and disagreement; a shared joint encoder, bidirectional cross-view attention, and temporal refinement then drive joint-frame gating with residual correction. We evaluate the framework quantitatively on a synthetic Unity dataset with avatar ground truth and the real-world Ski-PTZ-Pose dataset with a constructed reference, and qualitatively on author-collected dual-view recordings. Relative to canonical averaging, CanonFuse3D reduces the mean per-joint position error by 44.5\% on Unity and 48.6\% on Ski-PTZ-Pose, while comparisons with learned and geometric references reveal metric- and input-dependent trade-offs rather than universal superiority. Controlled studies characterize the useful range across camera separations, small timing and sampling-rate errors, and pose or image occlusion, including failures under severe corruption. The same architecture supports alternative pose estimators through lightweight adaptation and additional cameras through pairwise composition.
\end{abstract}

\begin{keyword}
3D human pose estimation \sep dual-view fusion \sep post-estimation pose fusion \sep calibration-free refinement \sep human-centric representation \sep front-end generalization
\end{keyword}

\end{frontmatter}

\linenumbers

\section{Introduction}

Three-dimensional (3D) human pose estimation from a single camera has advanced rapidly, from temporal lifting of 2D keypoints~\cite{pavllo2019videopose3d,zheng2021poseformer,li2022mhformer} and pretrained motion representations~\cite{zhu2023motionbert} to video-based body models~\cite{kocabas2020vibe} and foundation-scale full-body estimators that operate on in-the-wild images~\cite{yang2026sam}. These estimators make 3D motion analysis practical for sports, where joint trajectories support performance evaluation, skill assessment, and injury prevention. Their outputs nevertheless inherit the limitations of a single viewpoint: depth along the optical axis is ambiguous, limbs are frequently self-occluded, errors depend on the body orientation relative to the camera, and per-frame predictions jitter over time. Skiing exposes all of these weaknesses at once, because athletes move rapidly, are observed from a distance, and change orientation continuously~\cite{rhodin2018learning,chen2021sportscap}.

A second camera offers a complementary observation of the same motion: a limb hidden in one view may remain visible in the other, and the two depth estimates need not share the same error. The practical question is how to exploit this additional evidence. Binocular reconstruction methods jointly estimate stereo keypoints or refine 2D and 3D poses using geometric constraints~\cite{wan2025rsb,wan2024dualdiffusion}; more general calibrated methods triangulate detections or exchange image features under known camera geometry~\cite{iskakov2019learnable,he2020epipolar,liao2024multiple,zhang2021adafuse}. Temporary placement and pan--tilt--zoom operation in skiing make such geometric information harder to maintain. Unknown-calibration methods relax this requirement: MetaPose, for example, initializes from existing monocular 3D estimates and refines pose and camera variables using 2D evidence~\cite{usman2022metapose}. Thus, the remaining distinction is not whether a method can reuse a monocular estimator, but what information its cross-view refinement requires.

We study the setting in which both cameras already have an off-the-shelf monocular estimator, and only their independently predicted 3D sequences are passed to a subsequent fusion stage. Given two approximately synchronized streams on a common joint set, the objective is to produce a more accurate and temporally stable canonical pose sequence without supplying images, 2D detections, camera intrinsics, or extrinsics to that stage. These restrictions concern the fusion module, not the upstream estimators that process RGB frames; the output is a human-centric pose representation rather than a recovered world-coordinate trajectory. This formulation asks how much a second predicted stream can improve the available poses when calibrated image-space reconstruction is outside the deployment interface. Additional cameras are considered separately through pairwise composition.

Figure~\ref{fig:teaser} summarizes the proposed CanonFuse3D framework. Both streams are first expressed in human-centric coordinates defined by body landmarks and aligned to each other by a sequence-level similarity transform, which exposes where the two estimates agree and where they conflict. For every joint and frame, the framework combines the view's own pose, the aligned opposite-view pose, their difference, and motion cues; a shared joint encoder, bidirectional cross-view attention, and short-range temporal refinement then produce a joint-frame reliability gate that weights the two view-specific pose candidates, together with a learned residual correction. The upstream estimator remains separate from the fusion model, so the framework can be placed behind any monocular 3D pose estimator that provides the common joint set.

\begin{figure}[t]
  \centering
  \includegraphics[width=0.98\linewidth]{figures/journal/study_overview.pdf}
  \caption{System overview. Two monocular 3D pose streams are expressed in human-centric coordinates, aligned in pose space, and fused using cross-view and temporal features. CanonFuse3D produces a refined canonical pose sequence without camera parameters.}
  \label{fig:teaser}
\end{figure}

We evaluate the framework on three skiing datasets: a synthetic Unity dataset with avatar ground truth, the public Ski-PTZ-Pose dataset with a constructed canonical reference, and author-collected indoor and outdoor recordings used for qualitative deployment. Beyond clean-input accuracy, controlled studies examine how the same model behaves when the camera separation, timing, or sampling rate changes, when pose evidence or image regions are missing, when the upstream estimator is replaced, and when more than two cameras are available.

The main contributions are:
\begin{itemize}
    \item A human-centric framework for calibration-free post-estimation dual-view 3D pose fusion, combining landmark canonicalization, sequence-level alignment, cross-view attention, temporal refinement, and reliability-gated residual fusion.
    \item Quantitative evaluation on two skiing datasets with explicit reference regimes, separating canonical-space comparisons from an independent native-output audit and distinguishing pose-only baselines, published-method adaptations, and additional-input geometric references, complemented by qualitative deployment examples.
    \item A systematic robustness analysis covering camera separation, timing and sampling-rate errors, and both pose-level and image-level occlusion, which identifies the useful operating range and its failure cases.
\end{itemize}

\section{Related Work}
\label{sec:related_work}

\subsection{Monocular 3D Human Pose Estimation}
Two-stage monocular pipelines first detect 2D joints and then lift their trajectories to 3D. VideoPose3D uses dilated temporal convolutions for this lifting step~\cite{pavllo2019videopose3d}; PoseFormer models spatial and temporal relations with transformers~\cite{zheng2021poseformer}, and MHFormer combines multiple pose hypotheses~\cite{li2022mhformer}. MotionBERT learns transferable motion representations and supports 3D recovery from 2D pose sequences~\cite{zhu2023motionbert}. Their input representation separates lifting from image detection, but errors or missing evidence in the 2D tracks can propagate into the reconstructed motion. Our alternative-front-end study isolates this lifting stage using matched 2D tracks.

Image- and video-based body estimators instead obtain pose and shape from visual input. SAM 3D Body estimates body, hand, and foot pose and shape from a single image~\cite{yang2026sam}; VIBE uses video features and an adversarial motion prior for temporally coherent body estimation~\cite{kocabas2020vibe}. Both categories can supply a monocular 3D stream, but neither category by itself reconciles predictions from a second synchronized camera. CanonFuse3D operates after that upstream inference and uses only the resulting common-joint 3D sequences.

In sports applications, the monocular problem is harder because athletic movements are rapid, atypical of common training corpora, and frequently self-occluded. Closest to our setting, Rhodin et al. introduced the Ski-Pose PTZ-Camera dataset of alpine skiers observed by rotating cameras and used multi-view consistency as supervision for a monocular model~\cite{rhodin2018learning}, and SportsCap targets monocular sports videos with sub-motion priors~\cite{chen2021sportscap}. Both remain monocular at inference; complementary 3D predictions from simultaneously available viewpoints are not fused after estimation.

\subsection{Dual-view and Multi-view 3D Pose Estimation}
Dedicated binocular methods address the loss of geometric redundancy when only two views are available. RSB-Pose jointly estimates stereo keypoints using disparity-indexed features and refines reconstructed poses with a pretrained pose transformer to handle short-baseline uncertainty and occlusion~\cite{wan2025rsb}. Dual-Diffusion jointly denoises binocular 2D and 3D poses, using geometric relationships and baseline-dependent normalization~\cite{wan2024dualdiffusion}. These methods exploit matched two-view observations during reconstruction. Their image/2D and stereo-geometry interfaces differ from a fusion module supplied only with independently estimated 3D streams.

A recent survey of markerless multi-view 3D human pose estimation emphasizes the benefits of complementary viewpoints while identifying occlusion, camera variation, limited supervision, temporal consistency, view selection, and the accuracy--complexity trade-off as continuing concerns~\cite{nogueira2025survey}. Classical and learned multi-view pipelines usually recover a single 3D pose inside the estimator. Learnable Triangulation combines per-view evidence through differentiable algebraic or volumetric triangulation~\cite{iskakov2019learnable}. Epipolar Transformers transfer image features along epipolar lines in a synchronized, calibrated setting~\cite{he2020epipolar}, and MVGFormer alternates appearance modules with learning-free geometric modules that triangulate refined 2D poses~\cite{liao2024multiple}. AdaFuse addresses unreliable or occluded views by estimating point correspondences between heatmaps and learning camera-specific fusion weights so that poor features do not corrupt stronger views~\cite{zhang2021adafuse}. Although these strategies use different fusion mechanisms, they operate on images, feature maps, heatmaps, or calibrated 2D observations as part of multi-view reconstruction rather than on independently produced 3D pose streams.

Other work relaxes calibration or supervision. MetaPose initializes pose and camera estimates from per-view monocular 3D predictions, then uses a learned optimizer with 2D keypoint uncertainty to refine them without supplied calibration or multi-view 3D supervision~\cite{usman2022metapose}. It is therefore closely related to post-estimation fusion, but its refinement retains 2D observation and camera-projection terms; our two-view adaptation is included in the benchmark with these additional inputs labeled. FLEX predicts view-invariant skeletal rotations and bone lengths from multiple streams without supplied extrinsics~\cite{gordon2022flex}. General multi-view formulations can also be applied to a pair of synchronized views, but reducing the camera count does not remove their observation or geometry requirements. The distinction in this study is the 3D-only refinement interface, not a claim that two-view fusion or reuse of monocular predictions is new.

Takahashi et al.~\cite{takahashi2018human} additionally recover temporal synchronization: using detected 2D joints as correspondences, they jointly estimate camera extrinsics, inter-camera time offsets, and 3D joint positions from unsynchronized, uncalibrated cameras. Synchronization is therefore recoverable from human motion itself. In this work the two input streams are assumed to be approximately synchronized, and the remaining timing tolerance is characterized experimentally.

\subsection{Pose-Space Refinement, Reliability, and Generalization}
Post-estimation refinement provides another way to improve predictions without replacing the upstream estimator. STRIDE is a model-agnostic test-time training method that fits a motion prior to each video and refines noisy initial 3D estimates into temporally coherent poses under occlusion~\cite{lal2025stride}. Its objective is single-video temporal refinement rather than the reconciliation of synchronized viewpoints. Reliability weighting also appears within multi-view estimators: Learnable Triangulation predicts confidence weights for algebraic triangulation, and AdaFuse estimates adaptive weights for camera-view features~\cite{iskakov2019learnable,zhang2021adafuse}. In contrast, the reliability gate in this study acts after monocular 3D estimation and selects between aligned 3D pose candidates at the joint--frame level. The distinction is important because upstream confidence, image-feature quality, and residual pose-space reliability are not interchangeable signals.

Robust deployment also depends on how an estimator responds to a changed pose or camera distribution. PoseAug learns differentiable changes to posture, body size, viewpoint, and position to diversify paired 2D--3D training poses and improve cross-dataset generalization~\cite{gong2021poseaug}. AdaptPose instead generates target-like 3D motions and camera orientations from source poses and uses target-domain 2D poses without requiring target 3D labels or camera information~\cite{gholami2022adaptpose}. Both approaches adapt the training distribution of a 3D pose estimator. Our experiments address a different transfer boundary: they hold the pose-space fusion architecture fixed, replace the upstream monocular estimator, and then compare frozen transfer with target-specific adaptation of the refinement model.

The resulting research question is whether two independently estimated, approximately synchronized monocular 3D streams can be refined into a more accurate and temporally stable canonical sequence when the fusion stage receives no images, 2D detections, or camera parameters. CanonFuse3D studies this restricted interface through human-centric pose alignment, cross-view temporal reasoning, and joint-frame gating. The supporting experiments characterize its benefits and failure conditions rather than treating it as a replacement for every image- or geometry-based reconstruction method.

\section{Proposed Method}
\label{sec:method}

\subsection{Overview}
\label{subsec:method_overview}

The fusion interface contains only the two 3D streams: images and 2D observations remain upstream, and neither supplied camera parameters nor a world-coordinate reconstruction is required or produced by this module.

CanonFuse3D takes independently estimated pose sequences \(P^L,P^R\in\mathbb{R}^{T\times J\times3}\) from two approximately synchronized views, where \(T\) is the number of frames and \(J\) the number of joints. It predicts a refined canonical pose sequence \(P^F\) of the same shape without camera parameters. Figure~\ref{fig:feature_construction} shows the complete pipeline. Stages 1--3 place both streams in human-centric coordinates and align them to each other in pose space; stage 4 assembles, for every joint and frame, the five three-dimensional cues that describe the view's own estimate, the aligned opposite estimate, their disagreement, local motion, and the alignment rotation; stages 5--8 encode these cues, exchange information across views and over time, and fuse the two view-specific pose candidates through a reliability gate and a residual correction.

\begin{figure}[!t]
  \centering
  \includegraphics[width=0.98\linewidth]{figures/journal/method_pipeline.pdf}
  \caption{CanonFuse3D pipeline and feature construction. Stages 1--3 canonicalize and align the two pose streams; stage 4 concatenates five three-dimensional cues (own-view pose, aligned opposite-view pose, their difference, temporal velocity, and alignment rotation) into a 15-channel feature at every joint and frame; stages 5--8 encode the features, apply bidirectional cross-view attention and temporal refinement, and fuse the two base candidates \(B^L\) and \(B^R\) through a joint-frame gate \(\alpha\) with a residual correction \(\Delta P\). The right stream uses the same construction with the view roles reversed.}
  \label{fig:feature_construction}
\end{figure}

\subsection{Human-Centric Landmark Canonicalization}
\label{subsec:canonicalization}

Independent monocular estimates can differ in global position and orientation. We express each stream relative to a fixed reference pose \(A\) using the hip midpoint, hip direction, and neck direction. Let \(h_L,h_R,n\) denote the left hip, right hip, and neck, and let \(\operatorname{nrm}(v)=v/(\|v\|_2+\epsilon)\). The origin and body axes are
\[
\begin{aligned}
c&=\tfrac12(A_{h_L}+A_{h_R}), &
e_x&=\operatorname{nrm}(A_{h_R}-A_{h_L}),\\
e_y&=\operatorname{nrm}(A_n-c), &
e_z&=\operatorname{nrm}(e_x\times e_y).
\end{aligned}
\]
An additional landmark-direction cue resolves the sign of \(e_z\). With the axis matrix \(R=[e_x,e_y,e_z]\), the transformed coordinates are
\[
\tilde{P}_{t,j}=(P_{t,j}-c)R.
\]
The left stream, right stream, and supervision reference use separate body references. Each origin and axis matrix remains fixed over the sequences to which it is applied; the representation is therefore not a per-frame root-relative pose. Skeleton scale is not normalized, and the constructed axes need not be orthogonal. In the evaluated protocol, the reference is shared within each batch. Its exact selection and landmark conventions are specified in Supplementary Section~\ref{S-sec:implementation_record}.

\subsection{Pose-Level Cross-View Alignment}
\label{subsec:pose_alignment}

Because the independently anchored streams can retain residual translation, orientation, and scale mismatch, we apply pose-level cross-view alignment in both directions. Specifically, the left pose sequence is aligned to the right pose sequence, and the right pose sequence is aligned to the left pose sequence:
\[
\tilde{P}^{L \rightarrow R} = \mathcal{A}(\tilde{P}^{L}, \tilde{P}^{R}),
\qquad
\tilde{P}^{R \rightarrow L} = \mathcal{A}(\tilde{P}^{R}, \tilde{P}^{L}),
\]
where \(\mathcal{A}(\cdot)\) denotes sequence-level pose alignment.

For each sample, let \(X,Y\in\mathbb{R}^{N\times3}\) contain the source and destination joint coordinates in row-vector form, with \(N=TJ\). We estimate one uniform scale \(s\), proper rotation \(Q\), and translation \(d\) for the whole sequence:
\[
(s^\star,Q^\star,d^\star)=
\underset{s\geq0,\;Q\in SO(3),\;d\in\mathbb{R}^{3}}{\arg\min}
\frac{1}{N}\left\|sXQ^\top+\mathbf{1}d^\top-Y\right\|_F^2.
\]
We solve this unweighted least-squares problem by singular value decomposition with the proper-rotation correction of Umeyama~\cite{umeyama1991least}, excluding reflections. Each direction uses one transform for all frames and joints in the sequence. These auxiliary pose-space transforms do not estimate camera extrinsics; they are also distinct from the per-frame reference alignment used only for evaluation.

The model requires dense, finite pose tensors and receives no joint-validity mask. Missing detections are represented by zero-filled raw poses, and the controlled pose-corruption study replaces or perturbs selected raw entries before canonicalization. All transformed entries participate equally in alignment and training, including corrupted or zero-filled entries; missing joints are neither interpolated nor omitted. Detection-failure flags are used only for reporting, and supervision references remain uncorrupted. Supplementary Section~\ref{S-sec:controlled_details} gives the intervention rules.

We further compute cross-view residuals:
\[
D^{L} = \tilde{P}^{R \rightarrow L} - \tilde{P}^{L},
\qquad
D^{R} = \tilde{P}^{L \rightarrow R} - \tilde{P}^{R}.
\]
These residuals encode disagreement between the two pose streams and provide cues about view-specific errors.

\subsection{Geometric-Motion Feature Construction}
\label{subsec:feature_construction}

As shown in stage 4 of Figure~\ref{fig:feature_construction}, for each view we construct joint-level geometric-motion features by concatenating pose, cross-view, and temporal cues. For the left-view stream, the feature is defined as
\[
X^{L}
=
\left[
\tilde{P}^{L},
\tilde{P}^{R \rightarrow L},
D^{L},
V^{L},
r^{R \rightarrow L}
\right],
\]
and for the right-view stream,
\[
X^{R}
=
\left[
\tilde{P}^{R},
\tilde{P}^{L \rightarrow R},
D^{R},
V^{R},
r^{L \rightarrow R}
\right].
\]
The velocity terms are first temporal differences of the canonical poses, \(V^{L}_{t}=\tilde{P}^{L}_{t}-\tilde{P}^{L}_{t-1}\) with \(V^{L}_{1}=0\), and likewise for \(V^{R}\). The terms \(r^{R \rightarrow L}\) and \(r^{L \rightarrow R}\) are the rotation vectors of the corresponding sequence-level alignment rotations, broadcast to all joints and frames of the sequence. Each cue is three-dimensional, so every joint at every frame is described by a 15-dimensional feature vector. The five cues give the network the view's own estimate, the opposite view expressed in the same pose space, the residual disagreement between them, local motion, and the global relation between the two streams; together they are the evidence from which the reliability of each candidate is inferred.

\subsection{Joint Feature Encoding}
\label{subsec:joint_feature_encoding}

The left and right geometric-motion features are passed through a shared joint-wise encoder:
\[
H^{L} = \phi(X^{L}),
\qquad
H^{R} = \phi(X^{R}),
\]
where \(\phi(\cdot)\) is a two-layer multilayer perceptron, each layer followed by layer normalization and GELU activation. The encoder maps the 15-dimensional input feature into a 128-dimensional hidden representation for each joint and frame.

Because the encoder is applied at the joint-feature level, it preserves the temporal and skeletal structure of the input sequence while producing a richer representation for subsequent cross-view interaction.

\subsection{Bidirectional Cross-View Attention}
\label{subsec:cross_view_attention}

To exchange complementary information between the two pose streams, we use bidirectional cross-view attention. For each frame, the left-view joint features attend to the right-view joint features, and the right-view joint features attend to the left-view joint features:
\[
\hat{H}^{L}
=
\mathrm{LN}\!\left(H^{L}+\mathrm{Attn}(Q=H^{L}, K=H^{R}, V=H^{R})\right),
\]
\[
\hat{H}^{R}
=
\mathrm{LN}\!\left(H^{R}+\mathrm{Attn}(Q=H^{R}, K=H^{L}, V=H^{L})\right).
\]
The attention is multi-head and is applied over the \(J\) joints within each frame, with separate parameters for the two directions. This allows the network to compare corresponding and related joint features across the two views. When a joint is unreliable in one view, the other view may provide complementary information. The residual connection and layer normalization stabilize the cross-view feature update.

\subsection{Temporal Refinement}
\label{subsec:temporal_refinement}

Frame-wise fusion alone may reduce spatial error but can still produce temporally unstable pose trajectories. To address this issue, we apply a temporal refinement module after cross-view attention.

The temporal refiner operates on each joint's hidden feature trajectory. The sequence is normalized and passed through a depthwise one-dimensional convolution of kernel size 5 along time, a GELU activation, dropout, and a pointwise convolution:
\[
\bar{H} = \hat{H} + \mathrm{TempConv}(\mathrm{LN}(\hat{H})).
\]
This residual formulation preserves the original cross-view representation while adding local temporal context. The temporal module is lightweight and targets short-range motion consistency, which is important for reducing jitter in sports pose sequences. The same module is applied to \(\hat{H}^{L}\) and \(\hat{H}^{R}\).

\subsection{Reliability-Gated Fusion and Residual Correction}
\label{subsec:reliability_fusion}

After cross-view attention and temporal refinement, the left and right hidden features are concatenated:
\[
Z = [\bar{H}^{L}, \bar{H}^{R}].
\]
A gate head predicts a joint-frame reliability gate:
\[
\alpha = \sigma(g_{\theta}(Z)),
\]
where \(g_{\theta}\) is a two-layer perceptron with GELU activation, \(\sigma\) is the logistic function, and \(\alpha \in [0,1]^{T \times J \times 1}\). The gate controls the relative contribution of the two view-specific pose candidates for each joint and frame.

The two streams have separate body references. The network combines their view-specific candidates through a learned pose-space mapping; this combination does not imply a shared physical coordinate frame.

We construct two base pose candidates:
\[
B^{L} = \frac{1}{2}(\tilde{P}^{L} + \tilde{P}^{R \rightarrow L}),
\qquad
B^{R} = \frac{1}{2}(\tilde{P}^{R} + \tilde{P}^{L \rightarrow R}).
\]
The reliability-gated base fusion is then computed as
\[
P^{base}
=
\alpha B^{L}
+
(1-\alpha)B^{R}.
\]

Finally, a residual head with the same two-layer structure predicts a pose correction:
\[
\Delta P = h_{\theta}(Z),
\]
and the final refined pose is
\[
P^{F} = P^{base} + \Delta P.
\]

The gate \(\alpha\) is interpreted as a learned reliability signal rather than a calibrated probabilistic uncertainty estimate. This distinction is important because the model is trained to improve fusion accuracy and temporal stability, not to produce calibrated uncertainty values.

\subsection{Training Objective}
\label{subsec:training_objective}

The model is trained with a supervised pose loss and four lightweight regularization terms:
\[
\mathcal{L}
=
\mathcal{L}_{sup}
+
\lambda_{view}\mathcal{L}_{view}
+
\lambda_{bal}\mathcal{L}_{bal}
+
\lambda_{ent}\mathcal{L}_{ent}
+
\lambda_{temp}\mathcal{L}_{temp}.
\]
The supervised term is the mean absolute error to the independently canonicalized reference \(\tilde{G}\),
\[
\mathcal{L}_{sup}=\frac{1}{3TJ}\sum_{t,j}\|P^{F}_{t,j}-\tilde{G}_{t,j}\|_{1}.
\]
The view-reconstruction term preserves the discrepancy between the two canonicalized inputs. With \(\delta=\tilde{P}^{L}-\tilde{P}^{R}\), the fused pose and gate reconstruct the two inputs as \(\hat{P}^{L}=P^{F}+(1-\alpha)\delta\) and \(\hat{P}^{R}=P^{F}-\alpha\delta\), and
\[
\mathcal{L}_{view}=\frac{1}{2}\left(\frac{1}{3TJ}\sum_{t,j}\|\hat{P}^{L}_{t,j}-\tilde{P}^{L}_{t,j}\|_{1}
+\frac{1}{3TJ}\sum_{t,j}\|\hat{P}^{R}_{t,j}-\tilde{P}^{R}_{t,j}\|_{1}\right).
\]
This term constrains the gate and fused estimate without treating the sequence-level similarity transforms as physical camera geometry. Two gate penalties avoid degenerate fusion: the balance term \(\mathcal{L}_{bal}=|\bar{\alpha}-0.5|\), where \(\bar{\alpha}\) is the mean gate over all joints and frames, discourages collapse to a single view, and the negative-entropy term
\[
\mathcal{L}_{ent}=\frac{1}{TJ}\sum_{t,j}\left[\alpha_{t,j}\log\alpha_{t,j}+(1-\alpha_{t,j})\log(1-\alpha_{t,j})\right]
\]
discourages overconfident gates during training. Finally, the temporal term penalizes second-order variation of the fused trajectory,
\[
\mathcal{L}_{temp}=\frac{1}{(T-2)J}\sum_{t=2}^{T-1}\sum_{j}\|P^{F}_{t+1,j}-2P^{F}_{t,j}+P^{F}_{t-1,j}\|_{2},
\]
encouraging stable motion without removing the supervised accuracy objective. The weights are \(\lambda_{view}=0.05\), \(\lambda_{bal}=0.02\), \(\lambda_{ent}=0.005\), and \(\lambda_{temp}=0.01\). All sums are additionally averaged over the mini-batch. No term uses an input-visibility mask to omit or reweight missing joints: an available reference still supervises the output at every joint, while a zero-filled input remains part of the view-reconstruction target.

\section{Experimental Protocol}

\subsection{Experimental Setup}

Each video view is processed independently by an upstream monocular pose estimator. CanonFuse3D and the pose-only baselines receive the resulting 3D sequences without images, 2D keypoints, or camera parameters. A separate geometric comparison group additionally receives estimated 2D joints and supplied camera calibration. All comparisons use matched sequences, joint conventions, and the stated human-centric coordinate protocol within each dataset.

\subsection{Data Sources and Annotation}
\label{subsec:dataset_annotation}

We use three skiing datasets with complementary roles: two for quantitative evaluation and one for qualitative deployment. Table~\ref{tab:dataset_summary} summarizes their content, scale, and reference regimes; Figure~\ref{fig:dataset_examples} shows paired-view examples.

\begin{table*}[!htb]
\centering
\caption{Data sources used in this study, with their content, scale, split, and reference regimes. Split counts are action-disjoint camera-pair sequences for Unity and camera-pair entries for Ski-PTZ-Pose; the author-collected videos are not used for training or quantitative evaluation.}
\label{tab:dataset_summary}
\footnotesize
\setlength{\tabcolsep}{3pt}
\renewcommand{\arraystretch}{1.08}
\begin{tabular}{@{}>{\raggedright\arraybackslash}p{0.125\linewidth}
                >{\raggedright\arraybackslash}p{0.18\linewidth}
                >{\raggedright\arraybackslash}p{0.18\linewidth}
                rrr
                >{\raggedright\arraybackslash}p{0.17\linewidth}@{}}
\toprule
Data source & Content & Scale & \multicolumn{3}{c}{Split} & Reference / role \\
\cmidrule(lr){4-6} & & & Train & Val & Test & \\
\midrule
Unity synthetic & Rendered dual-view videos, monocular 3D pose streams, and the avatar skeleton & 2 avatars \(\times\) 12 actions (24 sequences of 20--65 frames); 180 synchronized views per action (5 layers \(\times\) 36 azimuths) & 193,320 & 128,880 & 64,440 & Avatar 3D skeleton, 15 joints; training, testing, and all controlled studies \\
\midrule
Ski-PTZ-Pose & Real giant-slalom videos with 2D/3D annotations and PTZ calibration metadata & 12 sequences; 6 PTZ cameras; about 10K images & 150 & -- & 30 & Constructed canonical reference from the annotations, 13 joints; real-world evaluation \\
\midrule
Author-collected videos & Indoor and outdoor dual-view recordings with estimated pose streams & 6 recordings (12 videos); 360--2,609 frames per view & -- & -- & -- & No metric 3D reference; qualitative deployment only \\
\bottomrule
\end{tabular}
\end{table*}

\begin{figure*}[t]
\centering
\includegraphics[width=\textwidth]{figures/journal/dataset_examples.pdf}
\caption{Representative synchronized dual-view examples from the three data sources. Columns show (a,b) Unity synthetic data, (c,d) Ski-PTZ-Pose, and (e,f) author-collected videos; rows show the two input views. Unity and Ski-PTZ-Pose support quantitative evaluation under the reference regimes stated in the text. The author-collected recordings have no independent metric 3D annotations and are used only for qualitative illustration. The Unity panels are foreground-cropped for legibility; the other panels retain their original field of view.}
\label{fig:dataset_examples}
\end{figure*}

\paragraph{Unity synthetic dataset}
We rendered a virtual skiing scene in Unity~\cite{juliani2020} with one female and one male avatar performing 12 actions each, including steady skiing, braking, falls, and transitions between motion states. Each action is observed from 180 synchronized viewpoints arranged in five layers, each with 36 azimuths at 10-degree intervals. An experimental sample comprises one camera pair; the model always receives two pose streams. The avatar skeleton supplies 3D ground truth independently of the monocular predictions.

\paragraph{Ski-Pose PTZ-Camera dataset}
The public Ski-Pose PTZ-Camera dataset, denoted Ski-PTZ-Pose, contains competitive alpine skiers performing giant slalom runs observed by six synchronized pan-tilt-zoom cameras~\cite{rhodin2018learning,ski_poseptz_dataset}. It provides 3D poses, projected 2D annotations, and per-frame camera calibration for about 10K images in 12 sequences. For this study, view-specific 3D annotations are aligned to a selected reference camera by similarity transforms fitted on an initial segment; the coordinate-wise median of the aligned annotations on the remaining frames defines the supervision and evaluation reference. Errors therefore measure agreement with this constructed canonical reference, rather than independent physical 3D accuracy. Neither camera calibration nor the annotation-derived transforms enter CanonFuse3D.

\paragraph{Author-collected qualitative videos}
Six indoor and outdoor dual-view skiing recordings provide a third dataset for qualitative deployment. A pose estimator processes each view independently. These videos have no independent metric 3D annotations and are excluded from training and quantitative evaluation; they illustrate the visual coherence of the refined motion relative to its input estimates and direct averaging.

\paragraph{Joint convention}
The two quantitative datasets are evaluated on different joint sets because their references provide different joints (Table~\ref{tab:joint_convention}). The Unity avatar rig and the native SAM3D front end both provide the 15-joint skeleton, including two eye landmarks that also resolve the facing direction during canonicalization, so Unity is evaluated on all 15 joints. The Ski-PTZ-Pose annotations follow an H36M-style skeleton without eye landmarks, and so do the alternative lifting front ends; the intersection with the 15-joint skeleton is the common 13-joint subset (indices 2--14). Ski-PTZ-Pose and all cross-front-end comparisons therefore use this subset, and the main Unity comparison reports it as well so that both datasets are read on the same joints. The earlier controlled native-input Unity studies use the full skeleton. The matched RGB-occlusion and calibration study instead reports common-13 errors after all-15 canonicalization, as specified in Section~\ref{subsec:matched_geometry}.

\input{tables/joint_convention.tex}

\subsection{Implementation Details}
\label{subsec:implementation_details}

Each camera-pair sample is uniformly resampled to 30 shared frames. The refinement model uses a shared two-layer joint encoder with hidden dimension 128, four-head bidirectional cross-view attention, a depthwise temporal convolution of kernel size 5, and two-layer gate and residual heads, and is trained with AdamW. Unity uses action-disjoint training, validation, and test splits; Ski-PTZ-Pose uses its predefined training and test partitions. Unity model selection uses validation error, while Ski-PTZ-Pose uses predeclared checkpoints because no independent validation partition is available. Test errors do not determine model selection. Complete architectures, optimization settings, split sizes, and checkpoint selection details are provided in Supplementary Section~\ref{S-sec:implementation_record}.

Pose-only comparisons include the two single views, canonical averaging, aligned averaging, a 3D quality-weighted average, a temporal convolutional fusion network (TCN) trained on the same data, and SmoothNet~\cite{zeng2022smoothnet} applied to canonical averaging. Aligned averaging combines sequence-level similarity estimation~\cite{umeyama1991least} with equal averaging. Temporal convolutions are an established sequence-modeling approach~\cite{bai2018empirical}; our TCN is a local three-layer, non-dilated convolutional control, not a reproduction of that published architecture. The quality-weighted average is likewise a study-specific heuristic. SmoothNet is an adaptation of the published temporal refinement method to the averaged pose stream, with its architecture and training settings documented in the supplement.

The separate geometric group includes classical direct linear triangulation (DLT)~\cite{hartley2004multiple}, reprojection-filtered DLT, and DLT with a learned residual correction. The latter two are study-specific extensions of DLT, not reproductions of published learned triangulation systems. Their additional estimated 2D observations and camera calibration are stated explicitly in the main comparison table.

\subsection{Evaluation Metrics}

Errors use the coordinate convention of the corresponding evaluation protocol, without conversion to millimeters. Table~\ref{tab:native_reference} and the controlled canonical-space studies apply the stated body-coordinate transformation independently to predictions and references. Table~\ref{tab:native_output_comparison} instead evaluates admissible native outputs before that transformation and removes only per-frame pelvis translation for position error; PA-MPJPE remains a separate evaluation-only shape alignment. The equations below use the canonical-space notation; the native audit substitutes its declared prediction and reference coordinates. Lower values indicate better agreement with the corresponding reference, and neither coordinate protocols nor dataset reference regimes are pooled.

\subsubsection{Mean per-joint position error}
Mean per-joint position error (MPJPE) averages Euclidean distances between predicted and reference joint positions:
\[
E_{\mathrm{MPJPE}}=\frac{1}{TJ}\sum_{t=1}^{T}\sum_{j=1}^{J}\|P^F_{t,j}-\tilde{G}_{t,j}\|_2.
\]
Dataset-level results aggregate all evaluated joint-frame points. No additional reference alignment is applied for this metric.

\subsubsection{Procrustes-aligned position error}
Procrustes-aligned MPJPE (PA-MPJPE) evaluates pose shape after fitting a least-squares similarity transform from prediction to reference independently in each frame:
\[
E_{\mathrm{PA}}=\frac{1}{TJ}\sum_{t=1}^{T}\sum_{j=1}^{J}\left\|s_t P^F_{t,j}Q_t^{\top}+d_t-\tilde{G}_{t,j}\right\|_2,
\]
where, for each frame \(t\),
\[
(s_t,Q_t,d_t)=\underset{s\geq0,\;Q\in SO(3),\;d\in\mathbb{R}^3}{\arg\min}\;\sum_{j=1}^{J}\left\|sP^F_{t,j}Q^{\top}+d-\tilde{G}_{t,j}\right\|_2^2 .
\]
The fit allows translation, a nonnegative uniform scale, and a proper rotation, with no reflection. Each reported joint subset is aligned separately. Collapsed predictions are mapped to the reference centroid and retained in the average. This reference fit is used only for PA-MPJPE and does not change model inputs, selection, or other metrics.

\subsubsection{Acceleration error}
Acceleration error measures agreement between second-order pose differences:
\[
E_{\mathrm{acc}}=\frac{1}{(T-2)J}\sum_{t=2}^{T-1}\sum_{j=1}^{J}
\|\Delta^2 P^F_{t,j}-\Delta^2\tilde{G}_{t,j}\|_2,
\]
where \(\Delta^2 P_{t,j}=P_{t+1,j}-2P_{t,j}+P_{t-1,j}\). It measures temporal agreement at the evaluated sampling interval, in coordinate units per sampled frame squared.

\subsection{Robustness and Generalization Protocols}
\label{subsec:controlled_protocol}

The robustness studies use the frozen Unity checkpoint and the complete held-out split; additional training is restricted to the explicitly identified gate ablations and front-end adaptation experiments. Camera separation is analyzed in six 30-degree bins of circular azimuth difference between the two cameras. Timing robustness uses seven right-stream offsets from \(-2\) to \(+2\) frames, including half-frame conditions, and seven sampling-rate errors from \(-2\%\) to \(+2\%\); an auxiliary diagnostic additionally estimates larger offsets of up to five frames. Pose-level corruption crosses three affected-view modes, three patterns, and five ratios, giving 45 conditions, and image-level occlusion applies the same view modes and patterns at two ratios to the RGB frames before the native pose estimator, giving 18 conditions; the two ratios parameterize different interventions and do not imply equal physical severity. Front-end generalization replaces SAM3D by three lifting estimators applied to the same ground-truth 2D tracks, isolating the lifting stage, followed by frozen transfer and target-specific or mixed adaptation. Finally, the multi-view extension composes the frozen dual-view model over all camera pairs of nested one- to four-camera groups. Full intervention and aggregation rules are given in Supplementary Section~\ref{S-sec:controlled_details}.

A separate matched-input study compares five frozen methods on clean images and random ratio-0.5/1.0 masks applied to the left, right, or both views. Paired estimated 2D and 3D observations come from the same SAM3D detections. The geometry-assisted methods additionally receive calibration, with right-camera extrinsic rotation perturbations of $0^\circ$, $1^\circ$, and $3^\circ$. All seven conditions retain the full test cohort, including detection failures (Section~\ref{subsec:matched_geometry}).

\section{Results}
\label{sec:results}

\subsection{Primary performance and cross-dataset evidence}

Table~\ref{tab:native_reference} compares CanonFuse3D with the pose-only fusion baselines, adapted published methods, and the separately labeled geometric group on the common 13-joint subset of both datasets after the shared body-coordinate normalization. On all 64,440 held-out Unity camera-pair sequences, CanonFuse3D reduces MPJPE from 0.2792 for canonical averaging to 0.1551 (44.5\%). Within this canonical-space protocol, the right row has larger MPJPE than the left row (0.4617 versus 0.1624). This gap cannot by itself be attributed to worse native front-end predictions because the streams pass through separate batch-referenced coordinate transforms; the native-output audit below evaluates the upstream predictions in declared frames. Aligning the right stream to the left before averaging reduces the canonical-space error to 0.1622, and the quality-weighted average brings no further gain. CanonFuse3D improves on the best of these non-learned methods in MPJPE (0.1551 versus 0.1622), PA-MPJPE (0.0637 versus 0.0667), and acceleration error (0.0297 versus 0.0308), so the learned fusion adds accuracy and temporal stability beyond what pose-space alignment alone provides. It is also more accurate than the temporal convolutional baseline trained on the same data (0.1610 MPJPE, 0.0690 PA-MPJPE, 0.0339 acceleration error). SmoothNet applied to canonical averaging gives the lowest acceleration error in the group (0.0113) but leaves the position error at 0.2925, so temporal smoothness and positional accuracy must be read separately. On all 15 Unity joints, the corresponding reduction over canonical averaging is 46.3\% (0.1541 versus 0.2868; Supplementary Table~\ref{S-tab:native_reference_all15}).

On the 30 fixed Ski-PTZ-Pose test sequences, the two views disagree far more (0.3552 and 0.4462), and pose-space alignment alone does not resolve the disagreement (0.3537). CanonFuse3D reduces MPJPE to 0.1794, a 48.6\% reduction over canonical averaging, with PA-MPJPE 0.1269 and acceleration error 0.0311. On this small real-world index the temporal convolutional baseline reaches lower MPJPE and PA-MPJPE (0.1581 and 0.1144) at a higher acceleration error (0.0411), and SmoothNet again gives the lowest acceleration error (0.0151) with a position error close to averaging (0.3227). These values measure agreement with the constructed annotation reference rather than independent physical 3D accuracy, and the 30-sequence index provides limited resolution between the two learned models.

The geometric group quantifies what calibration buys. With the estimated 2D joints and the supplied calibration, plain DLT is less accurate than CanonFuse3D on both datasets (0.1983 and 0.3984), reprojection gating does not improve it, and only a learned residual correction on top of DLT reaches lower MPJPE (0.0964 and 0.1643), at the cost of requiring 2D observations and camera calibration that the pose-only setting excludes. This group is therefore not a common-input comparison with CanonFuse3D. On Ski-PTZ-Pose the geometric predictions are first centered at their own per-frame pelvis, consistent with the root-relative annotations; this preprocessing fits nothing to the test reference.

The published adaptations give a metric-dependent comparison under their stated contracts. Avg.~+~STRIDE~\cite{lal2025stride} remains close to canonical averaging, with Unity MPJPE 0.2799, PA-MPJPE 0.1174, and acceleration error 0.0316; its corresponding Ski-PTZ-Pose values are 0.3483, 0.1861, and 0.0387. The official-protocol DeciWatch adaptation~\cite{zeng2022deciwatch} reaches 0.1134, 0.0695, and 0.0132 on Unity and 0.2239, 0.1549, and 0.0168 on Ski-PTZ-Pose. It has lower Unity MPJPE and acceleration error than CanonFuse3D but higher PA-MPJPE; on Ski-PTZ-Pose, CanonFuse3D has lower MPJPE and PA-MPJPE while DeciWatch has lower acceleration error. MetaPose~\cite{usman2022metapose}, which additionally uses estimated 2D joints, reaches 0.2655, 0.1089, and 0.0400 on Unity and 0.4253, 0.1425, and 0.0634 on Ski-PTZ-Pose. These adapted methods differ in input, supervision, and temporal-recovery task, so the table reports their outcomes without treating the lowest cell in each column as a common-input winner.

\input{tables/native_reference.tex}
\input{tables/native_output_comparison.tex}

Table~\ref{tab:native_output_comparison} instead evaluates outputs before the shared body-coordinate normalization. On Unity, the pooled native SAM3D streams have MPJPE/PA-MPJPE/acceleration error 0.1019/0.0732/0.0369. The native left stream is slightly more accurate than the right (0.0967 versus 0.1070 MPJPE), while sequence alignment and quality weighting reduce PA-MPJPE to 0.0666 and 0.0660. Among the published-method adaptations, per-view STRIDE~\cite{lal2025stride} reaches 0.1072/0.0776/0.0357, per-view DeciWatch~\cite{zeng2022deciwatch} reaches 0.6315/0.3033/0.0119, and MetaPose~\cite{usman2022metapose} reaches 0.2171/0.1088/0.0382. DeciWatch therefore smooths its per-view inputs but incurs large position error under this native-view application; this row is distinct from the canonical-average recovery protocol in Table~\ref{tab:native_reference}.

Every method in Table~\ref{tab:native_output_comparison} is scored directly in its declared output frame without the CanonFuse3D body-coordinate transform. Monocular and temporal-refinement results pool independently scored left and right camera streams; MetaPose remains in its restored left-camera frame, and calibrated geometry remains in its supplied world frame. For Ski-PTZ-Pose, the constructed annotation reference has no independently supplied mapping into these frames, so all native-space values are withheld rather than created by fitting a transform to test targets. The literal unaligned average is likewise retained only as a coordinate-mixing diagnostic. CanonFuse3D appears only in the canonical-space table; no native score is inferred through an unverified inverse. Because the protocols differ by whether the shared normalization is applied, methods must not be ranked across the two tables.

\FloatBarrier
\subsection{Reliability-gate behavior}

Two experiments distinguish the contribution of trainable gating from the behavior of an already trained checkpoint. First, Table~\ref{tab:gate_retraining} crosses dynamic versus fixed-0.5 gating with the presence versus absence of an output residual. The three learned configurations are retrained independently for 100 epochs on the same Unity partition and seed 42, with identical optimization and validation selection rules. The fourth configuration is deterministic and has no trainable output path. All four use the same complete test index and evaluation batch size 256; Supplementary Section~\ref{S-sec:gate_audit} details the protocol. This new full-model control is distinct from the archived model in Table~\ref{tab:native_reference}.

\input{tables/gate_retraining.tex}

With the output residual retained, learned gating lowers MPJPE from 0.17279 to 0.14639 (15.3\%) relative to fixed weights. With dynamic gating retained, the output residual lowers MPJPE from 0.16334 to 0.14639 (10.4\%). The full model also has the lowest PA-MPJPE. These single-run comparisons support complementary contributions to position accuracy under this protocol, not a claim of seed-independent gains. Removing the residual slightly lowers acceleration error, and the deterministic configuration is smoothest but has the largest position error; temporal smoothness alone does not establish reconstruction quality.

Second, we inspect the frozen main-benchmark checkpoint. Each gated candidate already contains information from both raw views after alignment, so \(\alpha=1\) or 0 selects a candidate rather than an original camera in isolation. On all 15 Unity joints, the correlation of \(\alpha\) with right-minus-left candidate error is 0.651, and the larger weight selects the lower-error candidate in 91.7\% of non-tied joint-frame comparisons. Table~\ref{tab:gate_region_summary} and Fig.~\ref{fig:gate_diagnostics_full} summarize this same complete-split audit. The mean left-candidate weight is 0.6804, indicating an asymmetric preference rather than balanced averaging. Its across-joint mean variation is small, while the joint-frame standard deviation is 0.1122. These diagnostics concern errors of the actual aligned candidates in evaluation coordinates; they neither establish probabilistic calibration nor directly measure raw-camera reliability.

\begin{table}[t]
\centering
\caption{Region-wise reliability-gate statistics on the complete Unity test split without masking.}
\label{tab:gate_region_summary}
\footnotesize
\setlength{\tabcolsep}{6pt}
\renewcommand{\arraystretch}{1.08}
\begin{tabular}{lccc}
\toprule
Region & Joints & \(\alpha\) & \(1-\alpha\) \\
\midrule
Upper limbs & 6 & \(0.6879 \pm 0.0132\) & 0.3121 \\
Lower limbs & 6 & \(0.6686 \pm 0.0103\) & 0.3314 \\
All joints & 15 & \(0.6804 \pm 0.0144\) & 0.3196 \\
\bottomrule
\end{tabular}
\vspace{0.4em}
\begin{minipage}{0.94\linewidth}
\footnotesize
\textit{Note.} Values report the mean and population across-joint standard deviation of per-joint mean \(\alpha\), not the standard deviation across all joint-frame weights. Upper limbs are the six shoulder/elbow/wrist joints and lower limbs the six hip/knee/ankle joints in the filtered skeleton. The all-joint group additionally includes both eyes and the neck. Balance and entropy regularization do not enforce a mean of 0.5 or probabilistic calibration.
\end{minipage}
\end{table}

\FloatBarrier
\begin{figure}[htbp]
\centering
\includegraphics[width=\linewidth]{figures/journal/gate_diagnostics_full.pdf}
\caption{Frozen-checkpoint gate diagnostics on the complete Unity split (all 15 joints). (a) Clean joint-frame weight distribution. (b) Per-joint means under clean inputs, complete raw-pose zeroing of either stream, and midpoint switches of the zeroed stream. The latter conditions also disturb canonicalization and alignment; they do not isolate gate reliability.}
\label{fig:gate_diagnostics_full}
\end{figure}

Freezing this checkpoint and replacing only the gate by 0.5 raises clean common-13 MPJPE from 0.1551 to 0.2104 on Unity and from 0.1794 to 0.1842 on Ski-PTZ-Pose. These inference interventions differ from retraining a model to compensate for a fixed gate. Under complete raw-pose zeroing and midpoint switches, the original checkpoint retains only a small advantage over the fixed-gate intervention, and on Unity it is worse than canonical averaging in every such stress condition. Turning off its output residual reduces Unity error under these extreme perturbations but worsens clean error and all tested Ski conditions. Complete zeroing also disrupts the batch anchor and Sim3 alignment, so these results do not establish automatic recovery of a missing view. All interventions, including unfavorable outcomes, are reported in Supplementary Table~\ref{S-tab:gate_frozen_interventions}; fixed-index frame--joint examples are shown in Supplementary Fig.~\ref{S-fig:gate_fixed_examples}.

\FloatBarrier
\subsection{Camera-pair separation}

The controlled studies in this and the following three subsections vary one capture condition at a time while keeping the checkpoint, held-out split, and evaluation path fixed. The camera-angle breakdown in Table~\ref{tab:view_angle} does not reveal a failure regime. CanonFuse3D MPJPE decreases from 0.1563 in the 0--30 degree bin to 0.1528 around 120--150 degrees, while the gain over canonical averaging stays between 45.9\% and 46.6\%. Wider azimuth separation therefore supplies slightly more complementary evidence in this synthetic setup, but the narrow range cautions against claiming a strong optimum. The bins contain 8,640--12,600 mutually exclusive test camera-pair sequences and together cover the complete 64,440-sequence test split.

\input{tables/view_angle_combined.tex}

To account for the four action-specific observations associated with each camera geometry, we collapse results by unordered camera pair before inference. Every angle bin retains a positive fusion gain under paired Wilcoxon testing (all Holm-adjusted \(p<0.001\); rank-biserial correlation 0.994--1.000). The across-bin Kruskal--Wallis test is statistically significant (\(p=0.025\)) but has a negligible effect size (\(\epsilon^2=0.00049\)). Of 15 Holm-adjusted post-hoc comparisons, only 90--120 versus 150--180 degrees remains significant (\(p=0.0064\)), again with a small rank-biserial effect (\(r_{rb}=0.053\)). Thus, the fusion benefit is reliable within every tested bin, whereas camera separation explains little of its magnitude. Detailed paired-test statistics are provided in Supplementary Table~\ref{S-tab:view_angle_significance}; Supplementary Fig.~\ref{S-fig:view_angle_extension} shows the corresponding absolute-error curves.

\FloatBarrier
\subsection{Temporal offset and sampling-rate drift}

Figure~\ref{fig:temporal_sampling} summarizes two synchronization stress tests as relative MPJPE changes from each method's unperturbed value. Supplementary Table~\ref{S-tab:temporal_sampling} provides the absolute errors and gate diagnostics. At zero perturbation, all-15 MPJPE is 0.1541. Across \(\pm2\) frame offsets, it ranges from 0.1537 to 0.1555; the worst case, a two-frame positive offset, is only 0.91\% above zero. Gate-error correlation stays between 0.646 and 0.650 and view-preference accuracy remains above 91.4\%. The modest variation suggests that sequence-level alignment and the five-frame temporal refiner tolerate short local disagreement, but this experiment does not establish robustness to long delays or unsynchronized action phases.

Clock-rate perturbations are even smaller. At \(+2\%\), MPJPE increases by 0.23\%; negative perturbations up to \(-2\%\) change MPJPE by less than 0.02\%. Acceleration error is not monotonic because temporal resampling smooths as well as misaligns trajectories. Consequently, MPJPE is the cleaner measure of rate-drift sensitivity in this controlled test.

\begin{figure}[!htb]
\centering
\includegraphics[width=0.9\linewidth]{figures/extension/temporal_sampling.pdf}
\caption{Relative MPJPE change under right-stream temporal offset and sampling-rate error on the complete held-out split. Each curve is normalized by its own method-specific unperturbed MPJPE; positive values indicate increased error. Dotted horizontal lines mark zero change.}
\label{fig:temporal_sampling}
\end{figure}

\FloatBarrier
\subsection{Missing-pose evidence}

Figure~\ref{fig:masking_extension} reports all 45 conditions of the complete held-out-split missing-pose study; Supplementary Table~\ref{S-tab:masking_extension} gives exact errors and gains at the representative ratios 0.5 and 1.0. Corrupted entries use a fixed mixture of zero replacement (35\%) and Gaussian perturbation (65\%, standard deviation 0.08 in raw input coordinate units, before canonicalization). For random masking, the ratio is the independent joint-frame probability; in the both-view condition, masks and corruption noise are sampled separately for the two streams. Distal-subset masking applies that probability only to the same eight filtered joint indices used by the executable masking protocol. Temporal masking selects the stated fraction of joints and corrupts one independently positioned contiguous 10-frame interval for each selected joint within every 30-frame sequence; its ratio therefore denotes selected joints rather than the global fraction of tensor entries.

Single-view corruption tests whether the complementary stream can rescue degraded evidence, whereas both-view corruption probes the method when redundancy is removed. The curves compare fused and canonical-average MPJPE, making explicit where learned fusion retains an advantage and where both methods fail together. At ratio 0.5, single-view temporal masking retains gains of 37.0\% for left corruption and 37.9\% for right corruption; distal masking retains 28.0\% and 31.5\%, respectively. The advantage narrows for unstructured random corruption. Most importantly, when random corruption affects both streams at ratio 0.5, fusion is 3.3\% worse than canonical averaging, whereas bilateral temporal and distal corruption still retain gains of 26.9\% and 12.5\%. Thus, the checkpoint can exploit a clean complementary stream and structured missingness, but dense independent corruption in both inputs lies outside its reliable regime.

\begin{figure}[!htb]
\centering
\includegraphics[width=0.98\linewidth]{figures/extension/masking_robustness.pdf}
\caption{Complete held-out-split masking robustness for left-, right-, and both-view corruption.}
\label{fig:masking_extension}
\end{figure}

\FloatBarrier
\subsection{Image-level occlusion through the pose front end}

Figure~\ref{fig:image_occlusion_examples} illustrates the three image-level occlusion patterns. Square patches filled with the image mean color are placed on the projected joints selected by the same random, distal, and temporal rules as the pose-level study, with a side length of 12\% of the projected person height; the masked frames are then processed by SAM3D and the frozen fusion model, so the resulting errors include missed detections and distorted poses produced by the front end. Supplementary Table~\ref{S-tab:image_occlusion} lists all 18 conditions.

\begin{figure}[!htb]
\centering
\includegraphics[width=0.98\linewidth]{figures/journal/image_occlusion_examples.pdf}
\caption{Image-level occlusion examples on one Unity frame. Rows show the two input views and columns the three patterns at ratio 0.5: random joint-frame selection, distal-subset selection, and temporal selection of contiguous 10-frame intervals. Red outlines mark the mean-color patches placed on the selected projected joints; the number of masked joints in each panel is stated in its label. The outlines are drawn for illustration only and are not present in the images passed to SAM3D.}
\label{fig:image_occlusion_examples}
\end{figure}

CanonFuse3D improves on canonical averaging in 15 of the 18 conditions. At random ratio 0.5, gains are 31.9\%, 34.3\%, and 20.9\% for left-, right-, and both-view occlusion, respectively. All twelve distal and temporal conditions retain positive gains, ranging from 14.6\% to 43.4\% for distal masks and from 42.3\% to 45.9\% for temporal masks. These results extend the controlled robustness evidence to errors produced by the image-to-pose front end.

Random occlusion at ratio 1.0 reverses the advantage in all three affected-view modes: fused MPJPE exceeds canonical averaging by 18.0\%, 9.2\%, and 28.4\% for left, right, and both views. In the bilateral case, MPJPE is 0.6311 versus 0.4914, and the pair-aligned SAM3D detection-failure rate is 36.36\%. Failure rates for the corresponding left- and right-view conditions are 40.51\% and 32.22\%. Failed detections remain in the evaluation as zero poses, so these results include complete upstream detection loss. The aggregate measurements do not isolate how much of the degradation comes from failed detections, errors in detected poses, or the fusion response. An unoccluded complementary view therefore does not guarantee a fusion advantage under this severe intervention.

Figure~\ref{fig:image_occlusion_extension} compares fused MPJPE for image-level and direct pose-stream interventions. Their matching ratios control different perturbations. For example, bilateral random corruption at ratio 0.5 retains a 20.9\% gain in the image-level experiment but gives a 3.3\% disadvantage in the pose-level experiment (Supplementary Tables~\ref{S-tab:image_occlusion} and~\ref{S-tab:masking_extension}). This difference shows that the two protocols probe distinct error distributions; it does not establish an ordering of physical occlusion severity. The condition counts describe one frozen checkpoint, rather than independent training repetitions.

\begin{figure}[!htb]
\centering
\includegraphics[width=0.98\linewidth]{figures/extension/image_occlusion_robustness.pdf}
\caption{Fused MPJPE under image-level occlusion propagated through SAM3D and direct pose-stream corruption, for left-, right-, and both-view interventions. Solid lines with filled markers denote image-level masks; dashed lines with open markers denote pose-stream corruption. Matching protocol ratios do not imply equal physical occlusion severity. Errors use dataset coordinate units.}
\label{fig:image_occlusion_extension}
\end{figure}

\FloatBarrier
\input{sections/c07_results.tex}

\subsection{Front-end generalization and adaptation}

CanonFuse3D is trained on SAM3D streams, but the framework itself is front-end agnostic: it consumes only 3D pose sequences on the common joint set and can be placed behind any monocular estimator that provides them. Table~\ref{tab:frontend_generalization} first evaluates the frozen SAM3D-trained checkpoint on three alternative lifting front ends without any retraining. MotionBERT retains a positive gain of 8.0\% over its own canonical average, whereas PoseFormer and VideoPose3D degrade by 18.0\% and 12.8\%. Canonicalization removes global position and orientation differences but not estimator-specific limb bias and temporal statistics, so a frozen model remains coupled to the error distribution it was trained on, and transfer to a new estimator should be validated rather than assumed.

\input{tables/frontend_generalization.tex}

\input{sections/p1_results.tex}

\FloatBarrier
\subsection{Qualitative evaluation}

\begin{figure}[!htb]
\centering
\includegraphics[width=0.98\textwidth]{figures/journal/realworld_qualitative.pdf}
\caption{Qualitative outputs on author-collected outdoor and indoor skiing videos. Panels (a), (b), (d), and (e) show the synchronized input views; panels (c) and (f) compare the two single-view predictions, canonical average, and CanonFuse3D output. These videos have no independent metric 3D annotations and are not used for training or quantitative evaluation. These examples were also included in the earlier conference manuscript.}
\label{fig:real_world_comparison}
\end{figure}

Figure~\ref{fig:real_world_comparison} illustrates deployment on the author-collected skiing videos, where left- and right-view poses are estimated independently and no camera calibration is used during inference. The proposed result remains more coherent than either single-view prediction or direct averaging, especially when the two views disagree around arms, knees, and feet. Because no independent 3D ground truth is available, this comparison is qualitative and is not included in any reported accuracy metric. The examples also show the main limitation: if both upstream predictions contain similar errors, the refinement can still inherit part of that error.

Qualitatively, the refined skeleton tends to preserve the more stable limb configuration while avoiding abrupt compromises caused by direct averaging. This is particularly visible in poses with strong torso rotation, occluded lower limbs, or large differences between left- and right-view predictions. For sports analysis, such behavior is useful because downstream measures such as joint-angle trajectories, limb extension timing, and motion smoothness are sensitive to frame-level outliers. At the same time, the examples should be read as a post-estimation refinement result rather than as full 3D reconstruction from images; the method improves the consistency of available pose streams but does not recover information that is absent from both inputs.

\FloatBarrier
\section{Discussion}
\label{sec:discussion}

In the canonical-space comparison, CanonFuse3D outperforms the non-learned pose-only baselines on both datasets in position, Procrustes-aligned, and acceleration error. Its margin over aligned averaging supports the joint benefit of the learned refinement module. The separate, matched Unity retraining ablation further supports contributions from both gating and output residual correction to position accuracy, while exposing a trade-off with acceleration error; this evidence remains single-seed and does not cover retraining on Ski-PTZ-Pose. Comparisons with learned methods are metric-dependent: the local temporal convolutional baseline reaches lower position error on Ski-PTZ-Pose, while SmoothNet attains lower acceleration error at a substantially higher position error. The DeciWatch adaptation has lower Unity MPJPE and acceleration error but higher PA-MPJPE than CanonFuse3D; on Ski-PTZ-Pose, CanonFuse3D has lower position and shape errors while DeciWatch remains smoother. STRIDE and MetaPose do not improve on CanonFuse3D in this canonical protocol, but their adaptation contracts and inputs differ. Within the calibrated group, a learned DLT residual correction reaches lower position error than CanonFuse3D, at the cost of additional estimated 2D observations and calibration. These outcomes identify different operating points, not a universal winner. The separate native-output audit must not be read as the same ranking. The earlier perturbation, transfer, and view-scaling studies compare against canonical averaging. The matched RGB study additionally compares geometric controls under shared corruption and varied calibration quality (Section~\ref{subsec:matched_geometry}); differing information budgets still prevent a common-input ranking.

The combined evidence supports a narrower claim than generic calibration-free multi-view reconstruction. CanonFuse3D is effective as a pose-space refinement layer when its two inputs are approximately synchronized, share the training front end's error distribution, and contain complementary evidence. It requires no camera parameters and retains a large advantage across the tested Unity angle bins. Pairwise composition extends the frozen dual-view model to larger camera groups, but it does not learn a joint representation over more than two views. Under the evaluated serial implementation, the number of forwards grows from one pair at two views to three and six pairs at three and four views, increasing latency from 2.02~ms to 6.02 and 12.06~ms. Because two views capture the largest accuracy gain over a single view while additional views provide smaller improvements at approximately threefold and sixfold latency, the dual-view configuration is the main accuracy--cost operating point. This timing conclusion is implementation-specific: batching or parallelizing independent pairs could reduce wall-clock latency without changing the pair count. Likewise, tolerance to small time and clock perturbations does not establish robustness to unknown long delays, dropped frames, or independently varying frame rates.

The front-end experiments delimit the generalization claim. The architecture is estimator-agnostic and, after 20 epochs of adaptation on training and validation data, improves over canonical averaging on every tested front end, whether adapted to one target or to a mixture; a single mixed model serves all four. What does not transfer is a frozen checkpoint: the fusion weights are coupled to the error distribution of the estimator they were trained on, so deploying CanonFuse3D behind a new estimator should include a short adaptation step, and the adaptation trades native-front-end retention against target recovery. The alternative-front-end experiments use ground-truth Unity 2D tracks to isolate 3D lifting; comparisons that include image-detection noise for those estimators remain necessary.

The image-level study broadens the native-SAM3D evidence to include synthetic RGB occlusion and upstream detection failures. Gains in 15 of 18 conditions support a useful operating range, while the three random ratio-1.0 failures show that an unoccluded second view alone is insufficient to ensure an improvement over averaging. The masking study uses Unity projections to place artificial patches and one frozen SAM3D front end. It does not establish robustness to natural occluders, weather, motion blur, or image-level errors from alternative estimators. Those settings require separate evaluation, and the observed association between detection failures and fusion degradation does not identify a causal mechanism.

Motion blur can affect the refinement through a different error process. During an exposure, fast limb motion or camera motion spreads image evidence along a trajectory, which can reduce joint-localization reliability or cause the front end to miss the person. The resulting 3D errors may persist across adjacent frames, altering both velocity features and cross-view disagreement. A clearer second view may provide useful evidence, but similar blur-induced errors in both streams can leave the gate with little reliable information to select. Errors in the hip or neck landmarks used for the batch anchor can additionally perturb the coordinate transform. These are expected mechanisms, not measured blur effects in this study. Temporal refinement may suppress isolated jitter, but neither temporal smoothing nor the reported patch-occlusion results establish robustness to exposure-induced motion blur.

The reference regimes constrain interpretation. Unity provides an avatar skeleton in dataset coordinates that is transformed using the same batch-referenced rule; it is not a metric motion-capture measurement. Ski-PTZ-Pose is evaluated against a transformed fused reference derived from view-specific annotations; its result measures agreement with that constructed reference rather than independent physical 3D accuracy. Raw-space and transformed-space errors are therefore not interchangeable, no skeleton-scale normalization is performed, and none of these values should be relabeled as millimeters.

The body reference is shared within each evaluated batch, so changing its composition can change the representation. The matched occlusion study also exposes collapse when reference landmarks are missing: effects can extend to other frames in the same batch (Supplementary Section~\ref{S-sec:c07_matched_geometry}). The results do not establish invariance to regrouping or single-sequence deployment, and severe-corruption errors cannot be attributed to the prediction module alone. A sample-independent variant would require retraining and evaluation; the exact reference convention is documented in Supplementary Section~\ref{S-sec:implementation_record}.

The reliability gate is informative but not calibrated uncertainty. Its mean is explicitly shaped by balance and entropy regularization, so high gate-error correlation and preference accuracy do not establish probabilistic calibration. Because the frozen-checkpoint stress tests reuse one trained checkpoint, they do not quantify variation across independent training runs. Reported sample- or group-bootstrap intervals describe resampling uncertainty conditional on the evaluated checkpoint and sampling units, not variation across training seeds or folds. Between-dataset uncertainty also remains unresolved. External sports datasets, metric 3D references, longer asynchronous sequences, and naturally occurring image degradation remain necessary directions.

\section{Conclusion}
\label{sec:conclusion}

CanonFuse3D refines two independently estimated, approximately synchronized monocular 3D streams without images, 2D observations, or camera parameters at the fusion stage. In the canonical-space benchmark, the complete framework reduces position error relative to canonical averaging by 44.5\% on Unity and 48.6\% on Ski-PTZ-Pose. Learned temporal and additional-input geometric methods retain advantages on particular metrics or datasets, so the result supports a restricted 3D-only refinement interface rather than universal superiority. Controlled studies show useful tolerance to the tested camera separations, small timing and sampling-rate errors, and many pose- and image-corruption conditions; severe random image occlusion can reverse the gain even when only one input view is masked. The matched-image comparison shows that accurate supplied calibration can favor learned geometric correction, while the tested larger calibration error favors pose-only fusion on clean and moderately masked inputs; severe masking remains a failure regime. The architecture supports alternative upstream estimators through adaptation and three or four cameras through pairwise composition, with two views offering the main accuracy--cost operating point under the evaluated serial implementation. These fixed-partition, single-run results motivate practical use subject to validation of the upstream estimator and capture conditions, while leaving cross-seed variability and broader real-world generalization unresolved.

\section*{Acknowledgments}
We thank the dataset providers and the open-source research community for making this work possible. This research, in part, used computational resources of Pegasus provided by the Multidisciplinary Cooperative Research Program in the Center for Computational Sciences, University of Tsukuba. This work used computational resources of Pegasus provided by University of Tsukuba through the HPCI System Research Project (Project ID: hp260146).

\section*{Declaration of competing interest}
\textbf{Author confirmation required before submission.} The authors must approve the final declaration wording.

\section*{CRediT authorship contribution statement}
\textbf{Author confirmation required before submission.} Contributions must be approved by all seven authors.

\section*{Data and code availability}
Ski-PTZ-Pose is available from the dataset provider~\cite{ski_poseptz_dataset}. The release status and repository URL for the Unity data, trained checkpoint, evaluation manifests, and evaluation scripts must be confirmed by the authors before submission. The author-collected indoor and outdoor videos are used only for qualitative illustration; their sharing conditions must also be confirmed before submission. The local reproducibility package records exact checkpoint and estimator hashes.

\section*{Ethics and generative-AI disclosure}
The quantitative study uses a synthetic dataset and a previously released research dataset. The author-collected indoor and outdoor videos are used solely for qualitative illustration and do not contribute training samples or quantitative measurements. The authors must confirm and report any applicable institutional-review and participant-consent requirements for these recordings before submission. Journal-required disclosure of generative-AI assistance must also be reviewed and approved by the authors before submission.

\setlength{\bibsep}{4pt}
\bibliographystyle{elsarticle-num}
\bibliography{references}

\end{document}
